\documentclass[11pt, a4paper, bibliography=totoc, listof=totoc]{scrartcl}

% --- PAKETE FÜR GRAFIKEN & FORMATIERUNG ---
\usepackage[utf8]{inputenc}
\usepackage[ngerman]{babel}
\usepackage[T1]{fontenc}
\usepackage{lmodern}
\usepackage{microtype} % Bessere Worttrennung und Ränder

% --- GEOMETRIE & LAYOUT ---
\usepackage[left=30mm, right=25mm, top=25mm, bottom=25mm]{geometry}
\usepackage{setspace}
\onehalfspacing % 1,5-facher Zeilenabstand (Pflicht bei vielen Hausarbeiten)

% --- MATHEMATISCHE FORMELN ---
\usepackage{amsmath}
\usepackage{amsfonts}
\usepackage{amssymb}

% --- TABELLEN & GRAFIKEN ---
\usepackage{booktabs}   % Schöne, wissenschaftliche Tabellenstriche
\usepackage{tabularx}   % Tabellen mit flexibler Spaltenbreite
\usepackage{graphicx}   % Einbinden von Grafiken (eure Plots aus dem Notebook)
\usepackage{float}      % Präzise Platzierung von Grafiken mit [H]
\usepackage[labelfont=bf]{caption} % Fettgedruckte Beschriftungen

% --- CODE-LISTINGS (Falls Codefragmente in den Text sollen) ---
\usepackage{listings}
\usepackage{xcolor}
\definecolor{codegreen}{rgb}{0,0.6,0}
\definecolor{codegray}{rgb}{0.5,0.5,0.5}
\definecolor{codepurple}{rgb}{0.58,0,0.82}
\definecolor{backcolour}{rgb}{0.95,0.95,0.92}

\lstdefinitions{
    style=mystyle,
    backgroundcolor=\color{backcolour},   
    commentstyle=\color{codegreen},
    keywordstyle=\color{magenta},
    numberstyle=\tiny\color{codegray},
    stringstyle=\color{codepurple},
    basicstyle=\ttfamily\footnotesize,
    breakatwhitespace=false,         
    breaklines=true,                 
    captionpos=b,                    
    keepspaces=true,                 
    numbers=left,                    
    numbersep=5pt,                  
    showspaces=false,                
    showstringspaces=false,
    showtabs=false,                  
    tabsize=2
}
\lstset{style=mystyle}

% --- VERWEISE & HYPERLINKS ---
\usepackage[colorlinks=true, linkcolor=black, citecolor=black, urlcolor=blue]{hyperref}

% =========================================================================
% DOCUMENT START
% =========================================================================
\begin{document}

\hypersetup{pageanchor=false}
\begin{titlepage}
    \vspace*{2cm}
    \begin{center}
        {\large Machine Learning \par} 
        \vspace{1.5cm}
        {\large Informatik \par} \\
        \vspace{2.5cm}
        {\huge \textbf{Evaluation von Q-Learning-Agenten im stochastischen Informationsspiel Blackjack} \par}
        \vspace{1.5cm}
        {\large Projektdokumentation \par}
        \vspace{0.5cm}
        {\large BHH \par}
    \end{center}

    \vfill

    \begin{flushleft}
        Abgabe: \\
        Hamburg, 16.06.2026
        
        \vspace{1cm}
        
        \textbf{Vorgelegt von:} \\
        \vspace{0.2cm}
        \begin{tabular}{@{}l l l l}
            Namen:          & Arvid Freese   & Linus Paliga   & Tom Markmann \\
            Matrikelnummer: & ______         & ______         & ______       \\
            Adresse:        & ___            & ___            & ___          \\
                            & _____ Hamburg  & _____ Hamburg  & _____ Hamburg\\
        \end{tabular}
    \end{flushleft}
\end{titlepage}


\thispagestyle{empty}
\newpage

% --- 1. ABSTRACT (Korrektur: Vor dem Inhaltsverzeichnis) ---
\begin{abstract}
\noindent Diese Projektdokumentation untersucht die Anwendbarkeit von tabellarischem Q-Learning im Blackjack-Szenario. Das Hauptaugenmerk liegt auf dem Vergleich zweier unterschiedlicher Architekturen: Einem \textit{Baseline-Agenten}, welcher ausschließlich auf der Standard-Zustandsrepräsentation operiert, und einem erweiterten \textit{Counting-Agenten}, welcher Kartenzählinformationen (\textit{Running Count}, \textit{True Count}) nativ in seinen Zustandsraum integriert. Die Agenten wurden über $100.000.000$ Episoden hinweg unter Verwendung von drei unterschiedlichen Seeds trainiert und anschließend evaluiert. Die Ergebnisse zeigen deutliche Unterschiede in der erlernten Aktionsverteilung sowie der resultierenden Win Rate.
\end{abstract}

\vspace{1.5cm}

% --- 2. INHALTSVERZEICHNIS ---
\tableofcontents
\newpage

\setcounter{page}{1} % Seitenzählung startet ab hier mit 1

% --- 3. EINLEITUNG ---
\section{Einleitung \& Fragestellung}
\label{sec:einleitung}
Hier folgt die Einleitung. Beschreibt kurz die Relevanz von Reinforcement Learning. 
\subsection{Zentrale Fragestellung}
\textit{Kann ein auf Q-Learning basierender Reinforcement-Learning-Agent eigenständig erfolgreiche Strategien beim Blackjack erlernen, und inwiefern verbessert die native Integration von Kartenzählinformationen den Spielerfolg im Vergleich zu einer Standard-Strategie?}


% --- 4. THEORETISCHER HINTERGRUND ---
\section{Theoretischer Hintergrund \& Simulationsumgebung}
\label{sec:hintergrund}
Erklärung der Spielregeln und eurer Simulations-Konfiguration aus dem Custom Environment.

\subsection{Eigenschaften der Simulationsumgebung}
Die Parameter eures Setups:
\begin{itemize}
    \item Anzahl der Kartendecks ($n = 6$)
    \item Schnitttiefe des Schlittens (Penetration = $0.75$)
    \item Dealer-Regel: \textit{Stand on Soft 17}
\end{itemize}

\subsection{Grundlagen des Tabular Q-Learning}
Die Aktualisierung der Q-Werte erfolgt nach der bekannten Bellman-Gleichung:
\begin{equation}
    Q(s, a) \leftarrow Q(s, a) + \alpha \left[ r + \gamma \max_{a'} Q(s', a') - Q(s, a) \right]
\end{equation}

% --- 5. AGENTEN-DESIGN ---
\section{Daten, Vorverarbeitung und Agentendesign}
\label{sec:agenten_design}

\subsection{Datenquelle und Trainingsdaten}

Anders als bei überwachten Lernverfahren existiert kein vorab
erhobener Datensatz. Sämtliche Übergänge $(s_t,a_t,r_{t+1},s_{t+1})$
entstehen online durch die Interaktion zwischen Agent und selbst
implementierter Simulationsumgebung. Die Trainingsdatenverteilung
hängt damit von der jeweils aktuellen $\epsilon$-greedy Policy ab. Zu
Beginn erzeugt die hohe Exploration vielfältige Übergänge; im
späteren Training dominiert zunehmend die gelernte Policy.

Die Verwendung einer Simulation erlaubt eine sehr große Zahl
kostengünstiger Episoden und vollständig bekannte Rewards. Sie
erzeugt zugleich ein Validitätsrisiko: Der Agent optimiert exakt die
implementierten Regeln und kann Modellvereinfachungen ausnutzen, die
in einem realen Blackjack-Spiel nicht existieren.

\subsection{Zustandsrepräsentation und Vorverarbeitung}

Die beiden Agenten unterscheiden sich ausschließlich durch ihren
State Encoder. Dadurch lässt sich der Nutzen der zusätzlichen
Kontextmerkmale isolierter untersuchen als bei gleichzeitig
verschiedenen Lernalgorithmen.

\paragraph{Baseline-Agent.}
Der Baseline-Zustand ist
\begin{equation}
  s_t^{B}=(p_t,d_t,a_t).
\end{equation}
Alle drei Werte werden als Ganzzahlen gespeichert. Die finale
Q-Tabelle enthält für jeden Trainingsseed genau 280 beobachtete Zustände.

\paragraph{Counting-Agent.}
Der erweiterte Zustand ist
\begin{equation}
  s_t^{C}=(p_t,d_t,a_t,\tilde r_t,\tilde c_t,\tilde m_t).
\end{equation}
Zur Begrenzung des Zustandsraums werden die Zusatzmerkmale
diskretisiert. Dabei bezeichnet $\operatorname{trunc}$ die im Code
verwendete Ganzzahlumwandlung, die Nachkommastellen gegen null abschneidet:
\begin{align}
  \tilde r_t &= \operatorname{trunc}(\operatorname{clip}(r_t,-20,20)),\\
  \tilde c_t &= \operatorname{trunc}(\operatorname{clip}(c_t,-10,10)),\\
  \tilde m_t &= \operatorname{clip}(\lfloor m_t/52\rfloor,0,6).
\end{align}
Eine klassische Normalisierung ist für die Q-Tabelle nicht
erforderlich, weil keine kontinuierliche Funktion approximiert wird.
Das Clipping verhindert seltene Extremwerte, führt aber zu
Informationsverlust. Zudem sind Running Count, True Count und
verbleibende Decks mathematisch voneinander abhängig. Die
Repräsentation enthält somit Redundanz und vergrößert die Tabelle stark.

Am Ende des finalen Trainings enthalten die drei Counting-Tabellen
zwischen 113.744 und 113.769 Zustände. Im Mittel sind dies 113.760
und damit rund 406-mal so viele wie bei der Baseline. Nicht besuchte
Zustände erhalten implizit Q-Werte von null. Bei Gleichstand wählt
\texttt{argmax} die Aktion Stand; diese technische Tie-Break-Regel
beeinflusst insbesondere unbekannte oder kaum trainierte Zustände.

\subsection{Modellarchitektur}

Es wird kein Deep-Learning-Modell eingesetzt. Die Architektur besteht
aus einer Python-\texttt{defaultdict}, die für jeden beobachteten
Zustand einen Vektor mit zwei Q-Werten speichert:
\begin{equation}
  Q(s)=[Q(s,\text{Stand}),Q(s,\text{Hit})].
\end{equation}
Neue Zustände werden mit $[0,0]$ initialisiert. Nach jedem Schritt
erfolgt genau ein Online-Update gemäß Gleichung~\ref{eq:qlearning}.
Es gibt weder Mini-Batches noch Experience Replay oder Target
Network. Der Vorteil dieser Architektur ist ihre Nachvollziehbarkeit;
ihr Nachteil ist die fehlende Generalisierung zwischen ähnlichen Zuständen.

\subsection{Hyperparameter und Auswahlprozess}

Vor dem finalen Lauf wurden kombinierte Hyperparameterkonfigurationen
in drei Stufen getestet. Die Parameter wurden nicht einzeln in einer
vollständigen Ablationsstudie variiert; die Ergebnisse zeigen daher
eine pragmatische Auswahl, aber kein globales Optimum.

\begin{table}[H]
  \centering
  \caption{Stufen der Hyperparametersuche}
  \label{tab:tuning_stages}
  \small
  \begin{tabularx}{\textwidth}{lrrrX}
    \toprule
    \textbf{Stufe} & \textbf{Konfigurationen} &
    \textbf{Trainingsseeds} & \textbf{Training} & \textbf{Greedy Evaluation}\\
    \midrule
    A & 16 & 1 & 3 Mio. Episoden & 150.000 je Checkpoint\\
    B & 8  & 2 & 10 Mio. Episoden & 500.000 je Checkpoint\\
    C & 4 Counting + Baseline & 3 & 20 Mio. Episoden & 1 Mio. je Checkpoint\\
    \bottomrule
  \end{tabularx}
\end{table}

Für den finalen Lauf wurde die Counting-Konfiguration
\texttt{c05\_long\_default} ausgewählt. Sie war in Stufe C mit einem
mittleren Reward von $-0{,}05259$ die beste Counting-Variante und
zeigte mit einer Seed-Standardabweichung von $0{,}00052$ die höchste
Stabilität der untersuchten Counting-Konfigurationen. Die Baseline
lag bei 20 Millionen Episoden mit $-0{,}04906$ dennoch vorn. Das
sprach dafür, dem größeren Counting-Zustandsraum erheblich mehr
Training zu geben.

\begin{table}[H]
  \centering
  \caption{Hyperparameter des finalen Trainings}
  \label{tab:hyperparameters}
  \begin{tabularx}{\textwidth}{lX}
    \toprule
    \textbf{Hyperparameter} & \textbf{Wert und Umsetzung}\\
    \midrule
    Lernrate $\alpha$ & $0{,}01$, während des gesamten Trainings konstant\\
    Diskontierungsfaktor $\gamma$ & $0{,}99$\\
    Initiales $\epsilon$ & $1{,}0$\\
    Finales $\epsilon$ & $0{,}10$\\
    $\epsilon$-Decay & Linear über 90\,\% der 100 Mio. Episoden,
    danach konstant\\
    Trainingsepisoden & 100 Mio. je Agentinstanz\\
    Checkpoint-Intervall & 10 Mio. Episoden\\
    Batch Size & Nicht anwendbar, da Online-Update nach jedem Schritt\\
    \bottomrule
  \end{tabularx}
\end{table}

Die finale Konfiguration wurde für beide Agententypen verwendet, um
den Modellvergleich nicht zusätzlich durch unterschiedliche
Lernparameter zu verzerren.


% --- 6. TRAINING & EXPERIMENTELLES SETUP ---
\section{Training \& Experimentelles Setup}
\label{sec:setup}
Dokumentation eurer Trainingsmethode. Zur Sicherung der statistischen Robustheit wurden für beide Agenten-Typen drei orthogonale Trainings-Seeds verwendet: $\text{SEEDS} = \{1, 42, 123\}$. Die Gesamtzahl der Trainingsepisoden betrug pro Seed $100.000.000$.


% --- 7. EVALUATION & MODELLVERGLEICH ---
\section{Evaluation \& Modellvergleich}
\label{sec:evaluation}
Die Evaluation erfolgte ohne Exploration ($\epsilon = 0$) über $1.000.000$ Testepisoden pro Seed.

\subsection{Metriken im Vergleich}
Hier fügt ihr die Werte aus eurem finalen DataFrame (\texttt{per\_seed\_metrics\_df}) ein.

\subsection{Visualisierung der Lernkurve}
% Beispielhafte Einbindung eures Plots aus dem Notebook (Ordner 'images' muss existieren)
% \begin{figure}[H]
%     \centering
%     \includegraphics[width=0.85\textwidth]{images/learning_curve.png}
%     \caption{Verlauf des mittleren Rewards und des TD-Errors über die Episoden}
%     \label{fig:lernkurve}
% \end{figure}



% --- 8. DISKUSSION & VERBESSERUNG ---
\section{Diskussion, Limitationen und Ausblick}
\label{sec:diskussion}

\subsection{Beantwortung der Forschungsfragen}

\paragraph{Kann Q-Learning Kartenzählinformationen nutzen?}
Ja. Die True-Count-Bucket-Analyse zeigt systematische Unterschiede in
der Aktionswahl. Besonders deutlich ist die um fast 4,8 Prozentpunkte
niedrigere Hit-Rate des Counting-Agenten bei einem True Count von
mindestens drei. Dass sich die relative Performance gegenüber der
Baseline über die Buckets verändert, spricht ebenfalls für eine
aktive Nutzung. Ein strenger Kausalnachweis würde jedoch eine
Ablation benötigen, bei der einzelne Zusatzmerkmale entfernt oder
permutiert werden.

\paragraph{Verbessern die erweiterten Features die Performance?}
Im finalen aggregierten Experiment ja, aber nur geringfügig. Der
Counting-Agent verbessert den durchschnittlichen Reward um
$0{,}001794$ und die Win Rate um 0,087 Prozentpunkte. Dieser Vorteil
tritt in allen drei Trainingsseed-Paaren auf, ist wegen der sehr
kleinen Zahl unabhängiger Trainingsläufe aber nicht robust
signifikant. Dem Gewinn steht eine rund 406-mal größere Q-Tabelle
gegenüber. Für tabellarisches Q-Learning ist das Verhältnis von
zusätzlicher Komplexität zu Performancegewinn ungünstig.

\paragraph{Wie viele Episoden werden bis zur Konvergenz benötigt?}
100 Millionen Episoden reichen aus, damit Counting zur Baseline
aufschließt, nicht aber für einen belastbaren Konvergenznachweis. Die
Checkpoint-Leistung schwankt weiterhin und die Q-Tabelle wächst noch.
Auch die konstante Lernrate und die verbleibende Exploration von
$\epsilon=0{,}1$ erfüllen keine hinreichenden Bedingungen, um
theoretische Konvergenz zu behaupten. Längere Läufe könnten weitere
Erkenntnisse liefern, sollten aber mit zusätzlichen Seeds und
Checkpoint-Evaluationen kombiniert werden.

\subsection{Herausforderungen}

\paragraph{Exploration und großer Zustandsraum.}
Die Baseline besucht ihre kleine Zustandsmenge schnell und kann
Q-Werte häufig aktualisieren. Counting verteilt dieselbe Zahl an
Episoden auf mehr als 100.000 Zustände. Viele Kombinationen treten
selten auf, wodurch einzelne Schätzungen verrauscht bleiben. Ein
längerer $\epsilon$-Decay hilft bei der Abdeckung, verzögert aber die
Ausnutzung guter Policies.

\paragraph{Stochastik und Evaluation.}
Millionen Testepisoden reduzieren das Monte-Carlo-Rauschen einer
festen Policy stark. Sie ersetzen jedoch keine unabhängigen
Trainingsläufe. Die Ergebnisse von Seed 123 demonstrieren, dass ein
Modell trotz identischer Hyperparameter anders enden kann. Künftige
Aussagen über Signifikanz sollten deshalb auf mehr Trainingsseeds
statt ausschließlich auf noch mehr Episoden pro Eval-Seed beruhen.

\paragraph{Rechen- und Speicheraufwand.}
Das finale Training dauerte trotz paralleler Ausführung rund 13
Stunden. Counting-Checkpoints sind ungefähr 6,3 MB groß,
Baseline-Checkpoints nur rund 32 KB. Noch längere Läufe erhöhen
Laufzeit und Artefaktmenge, ohne dass tabellarisches Q-Learning
zwischen ähnlichen Zuständen Wissen übertragen kann.

\subsection{Grenzen der Aussagekraft}

Die Simulationsumgebung bildet nur einen Ausschnitt realer
Blackjack-Regeln ab. Naturals erhalten keinen 3:2-Payout; Double
Down, Split, Insurance, Surrender und variable Einsätze fehlen.
Gerade Kartenzählen wird in der Praxis überwiegend für die Anpassung
von Einsätzen verwendet. Der hier gemessene Nutzen unterschätzt oder
verändert daher potenzielle Effekte eines vollständigen Spiels.

Running Count, True Count und verbleibende Decks sind redundant. Dies
erleichtert dem Agenten möglicherweise bestimmte Entscheidungen,
vergrößert aber den Zustandsraum unnötig. Außerdem wird der True
Count ganzzahlig abgeschnitten und geclippt. Die Bucket-Evaluation
ordnet eine Episode nach ihrem Anfangszustand ein, obwohl sich der
Count während der Runde verändern kann. Die Aktionsraten
berücksichtigen zwar den aktuellen Count, der Episodenreward bleibt
jedoch dem Start-Bucket zugeordnet.

Die vorhandene Basic-Strategy-Visualisierung dient nur als
qualitative Hilfe. Da die Projektumgebung eine reduzierte Regelmenge
verwendet und die Visualisierung nicht alle Zustandsdetails des
Counting-Agenten realistisch rekonstruiert, ist sie kein formaler
Optimalitätsnachweis.

Auch die exakte Reproduzierbarkeit eines vollständigen Neutrainings
ist eingeschränkt. Umgebung und Action Space werden explizit
geseedet, die globale NumPy-Zufallsquelle für die $\epsilon$-greedy
Aktionswahl jedoch nicht pro Trainingsprozess gesetzt. Die
gespeicherten Modelle erlauben eine reproduzierbare greedy Evaluation
der berichteten Ergebnisse; ein erneuter Trainingslauf mit denselben
nominellen Seeds muss dagegen nicht bitidentisch enden. Künftig
sollte jeder Worker einen eigenen, explizit gesetzten
Zufallszahlengenerator verwenden.

\subsection{Verbesserungsmöglichkeiten}

Die unmittelbarsten Verbesserungen betreffen die Validität und
Effizienz des bestehenden Q-Learning-Experiments. Separate Agenten
mit ausschließlich True Count, Running Count oder verbleibenden Decks
würden in einer Ablationsstudie zeigen, welche Merkmale tatsächlich
Nutzen stiften. Auf dieser Grundlage könnte eine kompaktere
Zustandsrepräsentation entwickelt werden, die redundante
Kombinationen vermeidet. Ein längerer Lauf mit bis zu 500 Millionen
Episoden ist erst danach sinnvoll und sollte durch feste Checkpoints
sowie mindestens zehn Trainingsläufe abgesichert werden.

Die Umgebung sollte um Natural-Payout, Double Down, Split, Surrender
und variable Einsätze erweitert werden. Erst dann lässt sich
untersuchen, ob ein positiver Count tatsächlich in einen positiven
erwarteten Gewinn oder eine geeignete Einsatzstrategie übersetzt
werden kann. Eine passende Reward-Funktion könnte statt eines
konstanten Rundenergebnisses den Nettoertrag inklusive Einsatzhöhe abbilden.

Zusätzlich sollte das Evaluationsprotokoll stärker auf
Trainingsunsicherheit ausgerichtet werden. Mehr unabhängige
Trainingsläufe sind aussagekräftiger als eine weitere Vergrößerung
der bereits umfangreichen Evaluation einzelner Modelle. Explizit
gesetzte Zufallszahlengeneratoren pro Worker und unveränderliche
Manifeste für alle ausgewerteten Artefakte würden die
Wiederholbarkeit verbessern.

\subsection{Ausblick}

Über diese direkten Verbesserungen hinaus bieten sich mehrere
methodische Erweiterungen an. Transfer Learning könnte die
Baseline-Policy als Initialisierung oder Lehrer für den
Counting-Agenten verwenden. Dadurch müsste die grundlegende
Hit/Stand-Strategie nicht erneut für jeden Count-Zustand gelernt werden.

Für größere Zustands- und Aktionsräume bietet sich anschließend
Funktionsapproximation an. Ein Deep Q-Network (DQN) ist aufgrund des
diskreten Aktionsraums eine naheliegende wertbasierte Alternative und
könnte Wissen zwischen ähnlichen Zuständen teilen. Proximal Policy
Optimization (PPO) würde als Policy-Gradient-Verfahren einen
methodisch anderen Vergleich ermöglichen und insbesondere bei
zusätzlichen Aktionen relevant werden. Beide Verfahren führen jedoch
neue Hyperparameter, höhere Rechenkosten und mögliche
Trainingsinstabilität ein. Ein fairer Vergleich muss deshalb dasselbe
Seed-, Checkpoint- und Evaluationsprotokoll wie das tabellarische
Q-Learning verwenden.

Multi-Agent Reinforcement Learning ist für die aktuelle Aufgabe
weniger vordringlich, weil der Dealer einer festen Regel folgt und
nicht lernt. Es könnte erst bei mehreren interagierenden Spielern
oder adaptiven Gegenspielern relevant werden. Ein Einsatz in realen
Spielsituationen wäre nur nach einer deutlich realistischeren
Umgebung und einer rechtlichen sowie ethischen Bewertung sinnvoll.


% --- 9. FAZIT ---
\section{Fazit}
\label{sec:fazit}
Kurzes Schlusswort.


\newpage

% --- 10. LITERATURVERZEICHNIS ---
% \bibliographystyle{alpha}
% \bibliography{references}

% --- 11. ANHANG (Optional, z.B. Code-Auszüge, zusätzliche Grafiken) ---
\appendix
\section{Anhang}
\label{sec:anhang}

\subsection{Zentrale Artefakte}

\begin{table}[H]
\centering
\caption{Relevante Projektartefakte für die Reproduktion}
\small
\begin{tabularx}{\textwidth}{lX}
\toprule
\textbf{Artefakt} & \textbf{Pfad im Repository}\\
\midrule
Experiment-Notebook & \texttt{notebooks/01\_experiment.ipynb}\\
Ausgeführter finaler Lauf & \texttt{models/runs/big\_run\_launcher\_20260608\_170138/}\\
Finale Modelle & \texttt{models/20260608\_170145\_*\_agent.pkl}\\
Checkpoints & \texttt{models/checkpoints/<agent-seed>/}\\
Finale Metriken & \texttt{models/evaluations/20260610\_164317/}\\
Projektgrafiken & \texttt{images/}\\
Zusätzliche Berichtsauswertung & \texttt{Doku/figures/}\\
\bottomrule
\end{tabularx}
\end{table}

\subsection{Ausgewählte Ergebnisse der Hyperparametersuche}

\begin{table}[H]
\centering
\caption{Stufe C der Hyperparametersuche bei 20 Millionen Episoden}
\small
\begin{tabular}{llrrr}
\toprule
\textbf{Agent} & \textbf{Konfiguration} & \textbf{Avg. Reward} & \textbf{Seed-Std.} & \textbf{Q-States}\\
\midrule
Baseline & \texttt{c00\_default} & $-0{,}04906$ & $0{,}00138$ & 280\\
Counting & \texttt{c05\_long\_default} & $-0{,}05259$ & $0{,}00052$ & 108.681\\
Counting & \texttt{c09\_lr01\_gamma098} & $-0{,}05353$ & $0{,}00101$ & 108.597\\
Counting & \texttt{c00\_default} & $-0{,}05442$ & $0{,}00138$ & 108.590\\
Counting & \texttt{c15\_long\_high\_lr} & $-0{,}05528$ & $0{,}00106$ & 108.479\\
\bottomrule
\end{tabular}
\end{table}

\subsection{Reproduktionshinweise}

Das Experiment-Notebook erzeugt die finalen Evaluations-, True-Count- und Policy-Grafiken unter \texttt{images/} aus den gespeicherten Modellen und Evaluationsdaten. Für kompakte Berichtsdarstellungen wurden die finalen Metriken und True-Count-Ergebnisse aus den gespeicherten CSV-Dateien zusammengefasst. Zusätzlich wurden die 60 Checkpoints mit jeweils 100.000 Episoden ausgewertet. Die resultierenden Zusammenfassungen liegen unter \texttt{Doku/figures/}. Der LaTeX-Bericht wird unter Windows mit \texttt{build-latex.ps1} kompiliert.


\end{document}